\chapter{A closure theorem for the covering paradigm}\label{ch:closure}

Every record in this thesis, and every record we know of in this problem, is
built the same way: choose $\N$ so that each small prime offset $q$ has
$\N-q$ divisible by some small prime fixed in advance. Chapter~\ref{ch:barriers}
showed that pushing this paradigm to sub-$100$ digits runs into two hard
problems. The obvious next question is whether the paradigm itself can be
escaped---whether some other mechanism could certify $\N-q$ composite more
cheaply than a congruence.

This chapter argues that it cannot, within a language precise enough to make the
claim checkable. The result is not a theorem in the sense of a completed proof
in a journal; it is a theorem in the sense of a stated claim, with a proof
skeleton whose cases are individually established, two named conditionals, and
an adversarial falsification campaign that failed to break it. We are explicit
throughout about which is which.

\section{The empirical precursor}\label{sec:wp7}

Before formalising anything, we ran five bounded probes into candidate
mechanisms, each with a predeclared quantitative gate. The metric is
\emph{compression}: offsets certified per nat of design entropy, benchmarked
against covering (which manages roughly $4\,400$ offsets per nat at $r=3$,
falling to about $0.4$ per nat at $r\approx1\,200$).

\begin{description}
\item[P1. Prime-power moduli.] Upgrading a modulus $p$ to $p^{2}$ costs an extra
$\log p$ nats and yields \emph{zero} new kills: the class $q\equiv a\pmod{p^{2}}$
is a subset of the class mod $p$ that is already killed (at $p=7$ we count
$227\subset1\,613$). Strictly dominated.

\item[P2. Polynomial value-set identities.] Take $\N=m^{k}-c$, so that
$\N-q=m^{k}-d^{k}$ is algebraically composite whenever $q=d^{k}-c$. The best
case is $k=2$: $c=398$ kills $113$ offsets below $10^{5}$. But confining $\N$ to
a $k$-th-power family surrenders $(1-1/k)\log B$ nats---$115$ nats at
$B=10^{100}$---and the value set below $Q$ has at most $O(Q^{1/2})$ members. The
mechanism cannot scale past a few hundred offsets \emph{at any price}.

\item[P3. Sierpi\'nski--Riesel exponent families.] For $\N_n=k\cdot2^{n}+c$ the
residue mod $p$ cycles with period $\mathrm{ord}_p(2)$. For any fixed candidate
the induced constraint is one congruence class: ordinary covering, no entropy
gain. (The exponent structure does correlate residues \emph{across} the family,
which is a legitimate search-throughput trick---noted for the engine, not for
the entropy budget.)

\item[P4. Norm forms and splitting conditions.] ``$\N-q$ is a norm in $K$'' does
not certify compositeness---rational primes with the right splitting behaviour
are norms---and any usable splitting condition on $q$ is a congruence condition
modulo discriminant-related moduli. Covering in Galois-theoretic clothing.

\item[P5. Residual bias.] Could $\N$ be chosen so that residual complements are
composite \emph{more often} than the boost predicts? The banked campaigns answer
directly: across three covers and $1.2\times10^{9}$ candidates,
$|\Ehat-E|\le0.1$. Any unexploited compositeness structure is worth at most
$0.1$ nats.
\end{description}

All five terminated at their gates. That is suggestive but not conclusive: five
probes is five probes. The rest of this chapter turns the suggestion into a
claim with a proof structure, a formal language, and an adversary.

\section{The formal setting}

\begin{definition}[Construction scheme]\label{def:scheme}
Fix a bound $B$ and target $Q$. A \emph{construction scheme} is a pair
$(F,C)$:
\begin{itemize}
\item a \emph{candidate family} $F$: a parameter set $\Theta\subseteq\mathbb Z^{m}$
with a poly-size description and an evaluation map
$\N:\Theta\to2\mathbb Z\cap[1,B)$ computable in $\mathrm{poly}(\log B)$ time. Its
\emph{design cost} is $\mathrm{cost}(F)=\log B-\log|\N(\Theta)|$, the nats of
search freedom surrendered;
\item a \emph{certificate map} $C$: on input $(\theta,q)$ with $q<Q$ prime, it
outputs $\bot$ or a divisor witness $D(\theta,q)$ with
$1<D<\N(\theta)-q$ and $D\mid\N(\theta)-q$, verifiable in $\mathrm{poly}(\log B)$
time.
\end{itemize}
The certificate is \emph{planted} if the divisibility holds as an algebraic
identity over the family (checkable symbolically), and \emph{extracted} if $C$
discovers the divisor per instance. Write $K(\theta)$ for the certified offsets
and $U(\theta)$ for the residual ones.
\end{definition}

The \emph{effective exponent} $E_{\text{eff}}(F,C)$ is the logarithm of the
expected work, in candidates, to find a $\theta$ all of whose residual
complements are composite. For the classical cover, $E_{\text{eff}}=E$ of
\eqref{eq:E}, calibrated to $\pm0.1$ nats in \S\ref{sec:calibration}.

\begin{theorem}[Divisor-Paradigm Closure; target statement]\label{thm:closure}
Let $(F,C)$ be a construction scheme whose certificate map is planted and
expressible in the schema language $L$ of \S\ref{sec:grammar}. Assume:
\begin{itemize}
\item[(H1)] Hardy--Littlewood-type heuristics for shifted primes: conditional on
all divisibility data, the primality indicators of distinct complements are
independent with the Mertens-adjusted rates;
\item[(H2)] no compositeness test with constant soundness on random rough inputs
runs in $o(\text{one modular exponentiation})$.
\end{itemize}
Then $E_{\text{eff}}(F,C)\ge E^{*}(Q,B)-o(1)$, where $E^{*}$ is the covering
frontier. Moreover, extracted certificates on any window-dense family of targets
are factoring-equivalent, and certificate-free detection is already performed at
its known-optimal cost by the standard engine.
\end{theorem}

Informally: \emph{within $L$, nothing beats covering; outside planted-$L$, a
scheme must either factor or test, and both are at known cost floors.}

\section{The three lemmas}

\subsection{Lemma 2 first: correlations are covering in disguise}

We take the lemmas out of order because the second is the one that is actually
proved, and it does a surprising amount of work.

\begin{proposition}\label{prop:shared}
Let $\N$ be even and $q_1\ne q_2$ odd primes below $Q$. If $d>1$ divides both
$\N-q_1$ and $\N-q_2$, then $d\mid q_1-q_2$, and hence every prime factor of $d$
is smaller than $Q$.
\end{proposition}
\begin{proof}
Subtract: $d$ divides $(\N-q_1)-(\N-q_2)=q_2-q_1$, and $0<|q_1-q_2|<Q$.
\end{proof}

Trivial to prove, and yet it settles an entire class of speculative mechanisms.
Any scheme that hopes to make the residual complements \emph{jointly} luckier---%
choosing $\N$ so that its leftovers conspire to be composite together---must act
through a divisor shared by two complements, and Proposition~\ref{prop:shared}
says every such divisor is supported on primes below $Q$. But divisibility by
small primes is exactly what the cover fixes and the sieve harvests. There is
nothing left to exploit.

\begin{lemma}[Correlation exhaustion]\label{lem:correlation}
Let $\mathcal F_{<Q}$ be the information generated by the events
$\{s\mid\N-q\}$ for primes $s,q<Q$. Then every two-complement divisor event is
$\mathcal F_{<Q}$-measurable; conditional on $\mathcal F_{<Q}$ (and on the
sieve's data up to its bound), (H1) makes the residual primality indicators
independent, so the success probability factorises as $\prod_q(1-p_q)$ and
$E_{\text{eff}}$ decomposes as design cost plus $\sum_q p_q$---which is the
covering objective. Empirically, $|\Ehat-E|\le0.1$ over $1.2\times10^{9}$
candidates bounds any violation.
\end{lemma}

\subsection{Lemma 1: entropy accounting for planted certificates}

\begin{lemma}[Entropy accounting]\label{lem:entropy}
Let $(F,C)$ be planted. For each mechanism define its \emph{certified sub-image
density} $\delta_{\text{eff}}$: the density, within a generic window of width
$Q$ at height $B$, of integers for which the scheme can exhibit its planted
witness using poly-range parameters. Then the amortised design cost per
certified offset is at least $\log(1/\delta_{\text{eff}})-O(\log\log B)$, and
for every family in $L$ this is at least the covering cost.
\end{lemma}

The proof is by cases, and the cases are where the content is. We treat them in
\S\ref{sec:cases} after fixing the language, because the arguments are
family-specific and quantitative rather than uniform.

\subsection{Lemma 3: the detector floor, as an open problem}

Hypothesis (H2) deserves to be stated on its own, because it is the one
genuinely open ingredient and it is interesting independently of Goldbach
deserts.

\begin{conjecture}[H2, sub-modexp compositeness detection]\label{conj:h2}
Let $n$ be uniformly random among integers in $[X,2X]$ with no prime factor
below $B_0=(\log X)^{O(1)}$. Does there exist a randomised algorithm deciding
compositeness of $n$ with soundness $\ge2/3$ in time
$o(M(\log X)\cdot\log X)$---asymptotically cheaper than one modular
exponentiation?
\end{conjecture}

Known art: trial division and batch-gcd methods~\cite{Bernstein} certify exactly
the small-divisor mass and have soundness $o(1)$ on rough inputs; Jacobi symbols
are equidistributed on rough composites; Fermat, Euler, Miller--Rabin and
Frobenius tests all cost $\Theta(M\log X)$. We are aware of no sub-modexp result
in either direction, and no lower bound in any realistic model.

The downstream consumer is quantified: a detector at cost $c$ and recall $\rho$
improves throughput by $1/(c+(1-\rho))$, so $c=0.1$, $\rho=0.9$ is a $5\times$
speedup on every search in this thesis---and, since record difficulty is
logarithmic in compute, converts directly into records. The same primitive
dominates GIMPS and PrimeGrid~\cite{PrimeGrid}, so a positive answer would be
felt well beyond this problem.

\section{The schema language \texorpdfstring{$L$}{L}}\label{sec:grammar}

To make Theorem~\ref{thm:closure} checkable we fix a concrete grammar of
certificate mechanisms. Each schema compiles to three things: an enumerator of
the offsets it certifies, a symbolic divisor witness, and a design-cost term.
\[
\begin{array}{ll}
L::= & \mathrm{Congruence}(d,a)\quad \N\equiv a\ (\mathrm{mod}\ d)\\
   |\ & \mathrm{PolyImage}(f\in\mathbb Z[m],\deg\le4)\quad \N-q=A(m)B(m)\\
   |\ & \mathrm{ExpFamily}(\N=j\cdot a^{n}+c)\quad\text{per-offset order conditions}\\
   |\ & \mathrm{MultiPoly}(A,B\in\mathbb Z[u_1..u_v],\deg\le2,\ v\le3)\\
   |\ & \mathrm{NormForm}(K,\ \mathrm{disc}\le10^{4})\quad \N-q=\mathrm{Norm}_K(\alpha)\\
   |\ & \mathrm{Cyclotomic}(\Phi_n,\ n\le200)\\
   |\ & \mathrm{Compose}(L,L)\quad\text{depth}\le2
\end{array}
\]
Two exclusions are deliberate. Mechanisms whose certificate requires a
primality decision to define membership are circular and excluded; and
per-instance divisor discovery is \emph{extraction}, handled by the second horn
of the theorem rather than by $L$.

Two thresholds make ``beats covering'' precise. A schema must certify more than
$\mathrm{CEIL}=3\sqrt{Q}\approx1\,048$ distinct prime offsets---placing it above
the $\sqrt{Q}$ ceiling that limits identity families---and it must do so at a
marginal cost below covering's. The right benchmark for the second condition is
subtle and we got it wrong at first (\S\ref{sec:referee}); the correct one is
covering's \emph{average}:
\[
\frac{437\ \text{nats}}{10\,607\ \text{offsets}}\;=\;0.041\ \text{nats per offset},
\]
not its marginal endgame cost of about $1.01$. Covering wins on average, not at
the margin, because one modulus spends $\log r$ nats and certifies
$\pi(Q)/(r-1)$ offsets \emph{at once}, whereas every known algebraic mechanism
buys $O(1)$ offsets per unit of entropy. Together with the design-entropy budget
$\log B\approx461$, the two conditions imply that a bulk mechanism needs
$\mathrm{cost}<\log B/\mathrm{CEIL}=0.44$ nats per offset.

\section{Closing the families}\label{sec:cases}

\subsection{Univariate identity families: a ceiling, not a price}

A degree-$d$ form takes at most $Q^{1/d}$ values below $Q$, so it can certify at
most that many offsets \emph{regardless of budget}: $349$ at $d=2$, $50$ at
$d=3$, $19$ at $d=4$, $7$ at $d=6$. Since $\mathrm{CEIL}=1\,048$, every
$\mathrm{PolyImage}$ with $d\ge2$ is closed at any price, and $d=1$ is ordinary
covering. Cyclotomic schemas fall inside the same bound because $\varphi(n)\ge2$
for every $n\ge3$. This is the P2 probe generalised, and it retires two families
outright.

\subsection{ExpFamily: supply, not price}

$\mathrm{ExpFamily}$ is the strongest non-covering family in the grammar and
needs a quantitative argument. For $\N=j a^{n}+c$, a prime $s$ certifies the
offset $q$ when $(q-c)j^{-1}$ lies in the subgroup $\langle a\rangle$ mod $s$,
and the resulting condition confines the exponent $n$ to one class modulo
$\mathrm{ord}_s(a)$: a cost of $\log\mathrm{ord}_s(a)$ nats \emph{per offset}.

Two facts then close it. First, a structural floor: $s\mid a^{d}-1$ forces
$a^{d}>s$, so $\mathrm{cost}=\log d>\log\log_a s$---cheap certificates require
small moduli, which is exactly where covering is cheapest. Exhaustively over all
$s<10^{5}$ the cheapest is $\log 5=1.609$ (at $a=7$, $s=2801$).

Second---and this is the binding constraint---cheap certificates are
\emph{scarce}. A modulus $s$ serves only the offsets landing in its order-$d$
subgroup, a fraction $d/s$ of them, so it supplies about $\pi(Q)d/s$ offsets:
$s=2801$ with $d=5$ supplies about twenty. Solving the resulting purchase
problem exactly---sort every $s<10^{5}$ by $\log\mathrm{ord}$, buy offsets
cheapest-first---gives, for every base $a\in\{2,3,5,7\}$:

\begin{observation}[ExpFamily is over budget]\label{obs:expfamily}
Reaching $1\,048$ certified offsets costs $1\,151$ nats, against a total design
budget of $\log B\approx461$ nats: $2.5\times$ over. The most that any
budget-feasible $\mathrm{ExpFamily}$ design can certify is $419$ offsets.
\end{observation}

\subsection{MultiPoly: clusters, isolation, and orbits}

Multivariate identities are the interesting case, and the one where we initially
expected the theorem to break. Products of two binary quadratic forms have
\emph{window-dense} value sets---about $9\,300$ of the $17\,983$ prime offsets
below $Q=200\,000$ are representable---which blows straight past the univariate
$\sqrt Q$ ceiling. If those representations could be exhibited, the family would
be a genuine threat.

\paragraph{The cluster bound.} Let $G=A\cdot B$ be homogeneous of total degree
$d$ in $v$ variables, with value-set density exponent $\alpha=\min(1,v/d)$
(so the number of distinct values below $X$ is about $X^{\alpha}$). A
first-moment count of $k$-element clusters of $G$-values inside a width-$Q$
window below $X$ gives $X^{\alpha}(X^{\alpha-1}Q)^{k-1}$, hence clusters exist
only up to
\[
k_{\max}(\alpha)\;=\;1+\Bigl\lfloor
\frac{\alpha\log X}{(1-\alpha)\log X-\log Q}\Bigr\rfloor ,
\]
finite whenever $\alpha<1-\log Q/\log X$. At $X=10^{200}$, $Q=2\times10^{5}$
that threshold is $\alpha^{*}=0.9735$, and every combination with $v<d$ has
$\alpha\le(d-1)/d\le0.9$:

\begin{table}[h]
\centering\small
\caption{Cluster ceilings. A scheme needs about $18\,000$ certified offsets per
candidate; every shared-variable identity family with $v<d$ delivers $O(1)$.}
\label{tab:kmax}
\begin{tabular}{lrr@{\qquad}lrr}
\toprule
$(d_1,d_2,v)$ & $\alpha$ & $k_{\max}$ & $(d_1,d_2,v)$ & $\alpha$ & $k_{\max}$\\
\midrule
$(2,2,2)$ & $0.500$ & $2$ & $(2,2,3)$ & $0.750$ & $4$\\
$(2,3,3)$ & $0.600$ & $2$ & $(2,3,4)$ & $0.800$ & $5$\\
$(2,4,4),(3,3,4)$ & $0.667$ & $3$ & $(2,4,5),(3,3,5)$ & $0.833$ & $6$\\
$(3,4,5)$ & $0.714$ & $3$ & any $v<d$ & $\le0.9$ & $\le13$\\
\bottomrule
\end{tabular}
\end{table}

We checked the density exponent empirically for the one in-grammar case with
$v\ge3$: for $(x^2+y^2+z^2)(x^2+2y^2+3z^2)$, counting distinct values in
width-$10^{4}$ windows at heights $10^{6},10^{8},10^{10}$ gives
$\hat\alpha=0.761$ against the predicted $0.750$.

\paragraph{Above the threshold: isolation.} For $v\ge d$ the value set is
window-dense and existence no longer obstructs, so the closure must come from
\emph{access}. Around any known solution $u_0$, the lattice points $\delta$
keeping $G(u_0+\delta)$ inside a width-$2Q$ window form an ellipsoidal slab with
principal extents $Q X^{-(d-1)/d}$ along the gradient and
$Q^{1/2}X^{-(d-2)/2d}$ along the Hessian directions. At $X=10^{200}$, $Q=2\times
10^{5}$, $d=4$ these are $10^{-145}$ and $10^{-47}$: every extent is far below
$1$, so the only lattice point is $\delta=0$. Window solutions are pairwise
isolated; dense existence is non-constructive for geometric reasons, not
complexity-theoretic ones.

\paragraph{Symmetry cannot help.} One might hope an infinite symmetry group
$\Gamma$ moves between solutions. Three exhaustive cases: if $\Gamma$ preserves
$G$, navigation conserves the value---one offset per orbit. If $\Gamma$ preserves
one factor, it conserves that factor's value, so the certificate always exhibits
the same divisor; if that divisor is small enough to enumerate, the scheme is
sieve-equivalent, and if it is balanced ($\approx X^{1/2}$), we are back to the
isolated sliver. If $\Gamma$ mixes generators from both factors, the reachable
$(\log A,\log B)$ pairs form a coarse lattice with about $(\log X)^{2}\approx
2\times10^{5}$ points against a target window of relative width $10^{-195}$.
Any remaining route must solve $G(u)=m$ for \emph{prescribed} $m$---but a
product-form representation of $m$ \emph{is} a nontrivial factorisation of $m$,
which is the theorem's extraction horn, not a new mechanism.

\subsection{Compose: empty, not expensive}

Finally, depth-2 composition. Our scorer treats a composition's parts as
independent and \emph{sums} their certified offsets---generous---and Compose
still peaked well below threshold. But the generosity conceals a stronger fact:
two nontrivial family constraints on the same $\N$ are contradictory at record
scale. Requiring $\N=m^{2}-c_1$ and $\N=m'^{2}-c_2$ forces
$(m-m')(m+m')=c_1-c_2$, so $m\le|c_1-c_2|$, whereas $\N\sim10^{200}$ needs
$m\sim10^{100}$. For ordinary shifts the solution sets are tiny---$c=398,4686$
admits five solutions, the largest with $m=1073$. The intersection is empty at
scale, so the only viable composition is family-plus-congruence: covering, plus
one family capped at $349$.

\section{The falsification engine}\label{sec:engine-falsify}

A closure claim invites refutation, so we built the refuter ourselves. The
engine compiles a schema to (enumerator, witness, cost), evaluates it against
the live cover context, and applies the filter: a schema \emph{passes} if it
certifies more than $\mathrm{CEIL}$ distinct offsets at below-benchmark cost,
with the design budget enforced.

\paragraph{Tier 0: calibration.} Before trusting the engine we required it to
reproduce the hand-derived numbers of \S\ref{sec:wp7} from primitives. It does:
the best $m^{2}-c$ family below $Q=10^{5}$ is $c=398$ with $113$ killable
offsets, matching P2 exactly, and the apparent discrepancy in the
nats-per-residual-kill figure resolved as a $B=10^{100}$ versus $B=10^{200}$
context difference. Calibration also \emph{corrected} the record: covering's
endgame is about $1.0$ nats per offset at $Q=122\,000$, not the $2.5$ quoted
earlier from a different target.

\paragraph{Tier 1: exhaustive.} We swept $13\,661$ schema instances across all
seven families. \textbf{Zero passes.} The first run reported sixteen, which
turned out to be a filter-semantics leak---congruence-containing schemas being
scored against the escape criterion, when the covering paradigm cannot by
definition escape itself. After hardening, the near-miss board is topped by
$\mathrm{ExpFamily}$ at $3.1$--$4.2$ nats per residual kill, about $3\times$
covering's endgame, exactly as P3 predicted.

\paragraph{Tier 2: generative.} Finally we handed the problem to an adversary:
an LLM-driven evolutionary search~\cite{AlphaEvolve} whose fitness function was
our scorer. The contract was designed to be unfoolable: candidate programs may
\emph{only} emit schema JSON, they run in a locked-down subprocess with no
network, grading happens in a fresh interpreter with the trusted scorer, and a
pass requires hand audit before it counts. A hostile candidate that simply
prints $\texttt{fitness=1.0}$ verifiably scores $0$.

Forty-one evolved candidates produced zero passes and explored only four
families, converging on the same $\mathrm{ExpFamily}$ optimum the exhaustive
tier had already flagged as the closest near-miss and then stalling about
$53\times$ short of the bar.

\section{The referee broke before the theorem did}\label{sec:referee}

The most instructive thing that happened in this chapter was a failure of ours,
not of the theorem.

Within about ten candidates the evolutionary search drove its reported fitness
from $0.27$ to $0.46$, entirely inside $\mathrm{ExpFamily}$, by pushing the mean
per-offset cost from $3.71$ down to $2.20$ nats. It looked like a genuine climb
toward falsification. It was not: the search had found an accounting hole in our
scorer. $\mathrm{ExpFamily}$ was credited with about $10\,416$ certifiable
offsets while being charged $\log\mathrm{ord}$ nats \emph{per offset} with the
bill levied only against the $867$ residual primes---so a design was credited
with bulk coverage costing $10\,416\times2.20\approx22\,900$ nats of design
entropy against a total budget of $461$. A $50\times$ overdraft.

The audit protocol caught it before any claim was made, and the fix
\emph{strengthened} the mathematics. Charging per-offset mechanisms for every
offset they claim, and capping coverage at the affordable fraction
$\log B/\mathrm{cost}$, produced three things we did not have before: the
design-entropy budget constraint; the realisation that the honest benchmark is
covering's average $0.041$ nats per offset rather than its marginal $1.01$; and
the supply-side purchase argument of Observation~\ref{obs:expfamily}, which
closes $\mathrm{ExpFamily}$ quantitatively rather than empirically.

\begin{remark}[Methodological finding]\label{rem:referee}
Adversarial search against a formal claim is most valuable as a
\emph{referee-hardening} device, and should be reported that way rather than as
a failed refutation. The properties that made this work are worth restating,
because they are cheap and general: candidates emit data, never scores; grading
runs in a separate trusted process; every quantity is re-derived from
primitives; and a ``pass'' is a hand-audited event rather than an automatic
conclusion. Under those conditions an adversary optimising against your
evaluator is a free audit of your evaluator.
\end{remark}

\section{Verdict}

\begin{table}[h]
\centering\small
\caption{Per-family closure status. $13\,661$ enumerated plus $41$ evolved
schema instances produced zero passes; each family closes for its own
reason.}\label{tab:closure}
\begin{tabular}{lll}
\toprule
Family & Closure & Kind\\
\midrule
$\mathrm{Congruence}$ & is the covering baseline; composite $d$ dominated (P1) & proved\\
$\mathrm{PolyImage}$ ($k\ge2$) & value set $\le Q^{1/k}$, i.e. $\le349$ offsets & proved, price-free\\
$\mathrm{Cyclotomic}$ & $\varphi(n)\ge2$, same $\sqrt Q$ ceiling & proved, price-free\\
$\mathrm{ExpFamily}$ & purchase plan $1\,151$ nats vs $461$; cap $419$ & proved, quantitative\\
$\mathrm{MultiPoly}$ ($v<d$) & cluster bound $k_{\max}\le13$ & first moment\\
$\mathrm{MultiPoly}$ ($v\ge d$) & isolation + orbit conservation; else extraction & draft rigour\\
$\mathrm{NormForm}$ & reduces to $\mathrm{MultiPoly}$ & proved\\
$\mathrm{Compose}$ (depth 2) & intersection empty at $10^{200}$ & proved\\
\bottomrule
\end{tabular}
\end{table}

What remains genuinely conjectural, stated plainly:
\begin{enumerate}
\item \textbf{(H1)}, Hardy--Littlewood independence---standard but unproved,
supported here by the $1.2\times10^{9}$-candidate calibration.
\item \textbf{(H2)}, the sub-modexp detection question of
Conjecture~\ref{conj:h2}. A positive answer breaks the \emph{cost floor}, not
the covering monopoly.
\item The $o(1)$ terms in Lemma~\ref{lem:entropy}, and the
value-equidistribution inputs to the cluster counting.
\item \textbf{Mechanisms outside $L$.} The language is concrete and deliberately
finite. A mechanism that certifies compositeness by something that is neither a
planted divisor identity nor per-instance divisor discovery is not excluded by
anything here. We know of none, and the two structural horns---planted implies
entropy accounting, extracted implies factoring---suggest why, but this is the
honest frontier of the result.
\end{enumerate}

The programme-level consequence is worth stating because it changes what should
be worked on. If covering is not merely the best known paradigm but provably the
only competitive one within $L$, then records are governed entirely by the two
levers of Chapter~\ref{ch:frontier}: the rounding gap, and throughput. That is
the subject of the final chapter.
